\SourceLeaf{278}
\par\indent The numerical values in our example (for the Briggsian logarithm, with \(q\) expressed in degrees) are
\[
\begin{aligned}
\log m={}&-0{,}0049796163\ 94(\frac{q}{100})^3\\
&-0{,}0016150307\ 6(\frac{q}{100})^4\\
&-0{,}0023973954(\frac{q}{100})^5\\
&-0{,}0125671(\frac{q}{100})^6
\end{aligned}
\]
\Article{10}
\indent In an extensive survey, where very many points must be transferred from the spheroid to the sphere or in the opposite direction, it will be preferable to compute an extensive table once and for all, instead of going back to the formulas each time. Such a table will be easier to use if its argument is the latitude on the sphere, \(Q+q\), instead of the latitude on the spheroid, since the passage from the former to the latter is required much more often than the reverse. Anyone experienced in computation will scarcely need to be told that constructing such a table requires calculating only a moderate number of entries directly. From these, the rest can be determined with the same precision and very little effort by a suitable interpolation method. The series given in Articles 8 and 9 will therefore be used for this purpose. This is precisely why it is advantageous to avoid choosing a round number for \(P\) and to choose one for \(Q\).

I append such a table at the end of this memoir. It is based on the standard value \(Q=52^0 40'\) (as in the example considered so far) and runs through twelve degrees, from \(46^0 40'\) to \(58^0 40'\), covering every value of the argument \(Q+q\) at one-minute intervals. It gives the corresponding value of \(P+p\) to five decimal places of a second, and the Briggsian logarithm of \(m\) to ten places, in units of the tenth decimal place. It also gives the value of \(-\tfrac{\mathrm{d}m}{2m\mathrm{d}q}\), expressed in seconds. The use of this last column will be explained below. I have given the table to so many decimal places so that it will suffice even for the most precise calculation of a trigonometric survey, namely for a
